\BookRecto
\begingroup
\renewcommand{\runningtitle}{Unit 1. Nouns, the article, and the copula}
\hyphenpenalty=10000
\exhyphenpenalty=10000
\emergencystretch=3em
\newcommand{\reserve}[1]{\par\penalty-100\begingroup\dimen0=\pagegoal\advance\dimen0 by-\pagetotal\ifdim\dimen0<#1\newpage\fi\endgroup}
\newcommand{\itempara}[2]{\par\addvspace{3pt}\begingroup\raggedright\noindent\hangindent=1.8em\hangafter=1\makebox[1.8em][l]{#1.}#2\par\endgroup}
% book-source-body:begin

\addfontfeatures{Ligatures=TeXOff}
\raggedright
\booktitle{\BookComponent{ch01-learner}{Unit 1. Nouns, the article, and the copula}UNIT 1}{Nouns, the article, and the copula}

\reserve{4\baselineskip}
% mdparagraph:P001:begin
\section*{1. Spelling and case}\BookHeading{nav-ch01-learner-0001}{subsection}{1. Spelling and case}{1}\BookLabel{bk-ch01-learner-624dd0eb97a3}
% mdparagraph:P001:end

% teaching: CH01-GRAM-ORTHOGRAPHY CH01-GRAM-GENDER CH01-GRAM-CASES
% mdparagraph:P002:begin
\textbf{Æ æ} is a single vowel letter. \textbf{Þ þ} (thorn) and \textbf{Ð ð} (eth) represent the sounds written \emph{th} in modern English; the choice between the letters does not by itself distinguish the sounds of \emph{thin} and \emph{this}. This unit follows Cook's use of \textbf{ð}. His hooked \textbf{ę} is an editorial form of \emph{e}, retained in such words as \textbf{męte} and \textbf{męnn}. A macron marks a long vowel or diphthong: \textbf{ā, ǣ, ē, ī, ō, ū, ȳ; ēa, ēo, īe}. Preserve the marks printed in each word; do not add them by analogy.[1]
% mdparagraph:P002:end

% mdparagraph:P003:begin
Nouns have masculine, feminine, or neuter gender, independently of the sex or nature of the thing named. Learn the gender with the noun. A noun's declension determines its case endings; its nominative ending alone does not identify either gender or declension.
% mdparagraph:P003:end

% mdtable:T01:begin
% course-table-style: 2026-10
% course-panel:T01-P01:begin
\par\begingroup
\fontsize{10.5}{13.2}\selectfont
\setlength{\LTleft}{0pt}\setlength{\LTright}{\fill}
\noindent
\begin{tabular}{@{}>{\raggedright\arraybackslash}p{27.369mm}>{\raggedright\arraybackslash}p{93.819mm}@{}}
\toprule
\textbf{Case} & \textbf{Uses introduced here} \\
\midrule
Nominative (N) & Subject; a noun identifying the subject after “be.” \\
Genitive (G) & A dependent expressing possession or another relation, often English “of.” \\
Dative (D) & Recipient; the noun after the selected prepositional constructions below. \\
Accusative (A) & Direct object. \\
\bottomrule
\end{tabular}
\par\endgroup
% course-panel:T01-P01:end
% mdtable:T01:end

% mdparagraph:P004:begin
Singular and plural forms may coincide, and one spelling may serve several cases. Its construction and the agreement of accompanying words determine the analysis. The tables give all four cases, even when their forms are identical.[2]
% mdparagraph:P004:end

\reserve{4\baselineskip}
% mdparagraph:P005:begin
\section*{2. The article}\BookHeading{nav-ch01-learner-0002}{subsection}{2. The article}{1}\BookLabel{bk-ch01-learner-f6bfdc1fc647}
% mdparagraph:P005:end

% teaching: CH01-GRAM-ARTICLE
% mdparagraph:P006:begin
The definite article means \textbf{the}. It agrees with its own noun in gender, number, and case. The plural forms are common to all three genders.[3]
% mdparagraph:P006:end

% mdtable:T02:begin
% course-table-style: 2026-10
% course-panel:T02-P01:begin
\par\begingroup
\fontsize{10.5}{13.2}\selectfont
\setlength{\LTleft}{0pt}\setlength{\LTright}{\fill}
\noindent
\begin{tabular}{@{}>{\raggedright\arraybackslash}p{22.0mm}>{\raggedright\arraybackslash}p{22.688mm}>{\raggedright\arraybackslash}p{22.688mm}>{\raggedright\arraybackslash}p{22.688mm}>{\raggedright\arraybackslash}p{22.688mm}@{}}
\toprule
\textbf{Case} & \textbf{Masculine singular} & \textbf{Neuter singular} & \textbf{Feminine singular} & \textbf{All genders plural} \\
\midrule
Nominative & se & ðæt & sēo & ðā \\
Genitive & ðæs & ðæs & ðǣre & ðāra \\
Dative & ðǣm & ðǣm & ðǣre & ðǣm \\
Accusative & ðone & ðæt & ðā & ðā \\
Instrumental: reference & ðȳ, ðon & ðȳ, ðon & as dative & as dative \\
\bottomrule
\end{tabular}
\par\endgroup
% course-panel:T02-P01:end
% mdtable:T02:end

% mdparagraph:P007:begin
Recognize \textbf{ðȳ} and \textbf{ðon} as instrumental singular forms; their uses are not required in composition here. Nouns have no separate instrumental forms in these paradigms. The source also permits \textbf{ðām} for \textbf{ðǣm} and \textbf{ðǣra} for \textbf{ðāra}; the tables and exercises select \textbf{ðǣm} and \textbf{ðāra}.
% mdparagraph:P007:end

\reserve{4\baselineskip}
% mdparagraph:P008:begin
\section*{3. Strong masculine nouns}\BookHeading{nav-ch01-learner-0003}{subsection}{3. Strong masculine nouns}{1}\BookLabel{bk-ch01-learner-246e855423d2}
% mdparagraph:P008:end

% teaching: CH01-PAR-MASC CH01-PAR-MASC-E
% mdparagraph:P009:begin
\textbf{Fisc} illustrates the common consonant pattern. In \textbf{męte}, the nominative \textbf{-e} is replaced by \textbf{-es, -as, -a, -um} where those endings occur; it is not part of a stem to which another full ending is added.[4]
% mdparagraph:P009:end

% mdtable:T03:begin
% course-table-style: 2026-10
% course-panel:T03-P01:begin
\par\begingroup
\fontsize{10.5}{13.2}\selectfont
\setlength{\LTleft}{0pt}\setlength{\LTright}{\fill}
\noindent
\begin{tabular}{@{}>{\raggedright\arraybackslash}p{22.0mm}>{\raggedright\arraybackslash}p{22.688mm}>{\raggedright\arraybackslash}p{22.688mm}>{\raggedright\arraybackslash}p{22.688mm}>{\raggedright\arraybackslash}p{22.688mm}@{}}
\toprule
\textbf{Case} & \textbf{fisc singular} & \textbf{fisc plural} & \textbf{męte singular} & \textbf{męte plural} \\
\midrule
Nominative & fisc & fiscas & męte & mętas \\
Genitive & fisces & fisca & mętes & męta \\
Dative & fisce & fiscum & męte & mętum \\
Accusative & fisc & fiscas & męte & mętas \\
\bottomrule
\end{tabular}
\par\endgroup
% course-panel:T03-P01:end
% mdtable:T03:end

% mdparagraph:P010:begin
The plural of \textbf{męte} is displayed for recognition of the pattern; this unit uses its singular “food.” \textbf{Heofon, gāst, wæstm} follow \textbf{fisc}. Keep both syllables of \textbf{heofon} before endings. The divine name \textbf{God} has the singular \textbf{God, Godes, Gode, God}; no plural of this proper-name sense is required.
% mdparagraph:P010:end

% teaching: CH01-PAR-MASC-AE-A
% mdparagraph:P011:begin
In \textbf{dæg} and \textbf{hwæl}, singular \textbf{æ} alternates with plural \textbf{a}. The change does not apply to \textbf{wæstm}, whose \textbf{æ} stands before a consonant group.
% mdparagraph:P011:end

% mdtable:T04:begin
% course-table-style: 2026-10
% course-panel:T04-P01:begin
\par\begingroup
\fontsize{10.5}{13.2}\selectfont
\setlength{\LTleft}{0pt}\setlength{\LTright}{\fill}
\noindent
\begin{tabular}{@{}>{\raggedright\arraybackslash}p{22.0mm}>{\raggedright\arraybackslash}p{22.688mm}>{\raggedright\arraybackslash}p{22.688mm}>{\raggedright\arraybackslash}p{22.688mm}>{\raggedright\arraybackslash}p{22.688mm}@{}}
\toprule
\textbf{Case} & \textbf{dæg singular} & \textbf{dæg plural} & \textbf{hwæl singular} & \textbf{hwæl plural} \\
\midrule
Nominative & dæg & dagas & hwæl & hwalas \\
Genitive & dæges & daga & hwæles & hwala \\
Dative & dæge & dagum & hwæle & hwalum \\
Accusative & dæg & dagas & hwæl & hwalas \\
\bottomrule
\end{tabular}
\par\endgroup
% course-panel:T04-P01:end
% mdtable:T04:end

% teaching: CH01-PAR-MASC-SYNCOPE
% mdparagraph:P012:begin
\textbf{Fugol} and \textbf{morgen} lose their unstressed middle vowel before endings in the selected forms. This loss, called \emph{syncope}, is lexical: use the diagnostic forms given for each noun rather than deleting vowels from every word of two syllables. Forms retaining the vowel also occur.[4]
% mdparagraph:P012:end

% mdtable:T05:begin
% course-table-style: 2026-10
% course-panel:T05-P01:begin
\par\begingroup
\fontsize{10.5}{13.2}\selectfont
\setlength{\LTleft}{0pt}\setlength{\LTright}{\fill}
\noindent
\begin{tabular}{@{}>{\raggedright\arraybackslash}p{22.0mm}>{\raggedright\arraybackslash}p{22.688mm}>{\raggedright\arraybackslash}p{22.688mm}>{\raggedright\arraybackslash}p{22.688mm}>{\raggedright\arraybackslash}p{22.688mm}@{}}
\toprule
\textbf{Case} & \textbf{fugol singular} & \textbf{fugol plural} & \textbf{morgen singular} & \textbf{morgen plural} \\
\midrule
Nominative & fugol & fuglas & morgen & morgnas \\
Genitive & fugles & fugla & morgnes & morgna \\
Dative & fugle & fuglum & morgne & morgnum \\
Accusative & fugol & fuglas & morgen & morgnas \\
\bottomrule
\end{tabular}
\par\endgroup
% course-panel:T05-P01:end
% mdtable:T05:end

\reserve{4\baselineskip}
% mdparagraph:P013:begin
\section*{4. Strong neuter nouns}\BookHeading{nav-ch01-learner-0004}{subsection}{4. Strong neuter nouns}{1}\BookLabel{bk-ch01-learner-737b70e12605}
% mdparagraph:P013:end

% teaching: CH01-GRAM-WEIGHT CH01-PAR-NEUT-WEIGHT
% mdparagraph:P014:begin
Neuter nominative and accusative are identical in each number. The genitive and dative endings match those of \textbf{fisc}. In the ordinary monosyllabic pattern, the nominative/\allowbreak{}accusative plural has \textbf{-u} after a short stem, but no ending after a heavy stem. A short vowel followed by one consonant gives the short pattern here; a long vowel or long diphthong, or a short vowel followed by two consonants, gives the heavy pattern. A doubled consonant counts as two. Thus \textbf{hof} has \textbf{hofu}, whereas \textbf{sǣd, gærs, cynn} have no added N/\allowbreak{}A plural ending. Vowel length and syllable weight are not the same thing.[5]
% mdparagraph:P014:end

% mdtable:T06:begin
% course-table-style: 2026-10
% course-panel:T06-P01:begin
\par\begingroup
\fontsize{10.5}{13.2}\selectfont
\setlength{\LTleft}{0pt}\setlength{\LTright}{\fill}
\noindent
\begin{tabular}{@{}>{\raggedright\arraybackslash}p{22.0mm}>{\raggedright\arraybackslash}p{22.688mm}>{\raggedright\arraybackslash}p{22.688mm}>{\raggedright\arraybackslash}p{22.688mm}>{\raggedright\arraybackslash}p{22.688mm}@{}}
\toprule
\textbf{Case} & \textbf{hof singular} & \textbf{hof plural} & \textbf{sǣd singular} & \textbf{sǣd plural} \\
\midrule
Nominative & hof & hofu & sǣd & sǣd \\
Genitive & hofes & hofa & sǣdes & sǣda \\
Dative & hofe & hofum & sǣde & sǣdum \\
Accusative & hof & hofu & sǣd & sǣd \\
\bottomrule
\end{tabular}
\par\endgroup
% course-panel:T06-P01:end
% mdtable:T06:end

% mdparagraph:P015:begin
The heavy pattern also covers \textbf{lēoht, līf, gēar, dēor, anginn, lyft, weorc}. In \textbf{anginn}, apply the rule to the stressed stem \textbf{-ginn}, with \textbf{nn}. Plural forms of mass or abstract nouns do not automatically have the same sense as the singular; the vocabulary below marks those whose plurals are only for recognizing the morphology.
% mdparagraph:P015:end

% teaching: CH01-PAR-NEUT-TREO
% mdparagraph:P016:begin
\textbf{Trēo} has \textbf{w} before every ending and the exceptional N/\allowbreak{}A plural \textbf{trēowu}, despite its long vowel. Bare singular \textbf{trēow} also occurs.
% mdparagraph:P016:end

% mdparagraph:P017:begin
\textbf{trēo}\par\nobreak
% mdparagraph:P017:end

% mdtable:T07:begin
% course-table-style: 2026-10
% course-panel:T07-P01:begin
\par\begingroup
\fontsize{10.5}{13.2}\selectfont
\setlength{\LTleft}{0pt}\setlength{\LTright}{\fill}
\noindent
\begin{tabular}{@{}>{\raggedright\arraybackslash}p{39.459mm}>{\raggedright\arraybackslash}p{39.459mm}>{\raggedright\arraybackslash}p{39.459mm}@{}}
\toprule
\textbf{Case} & \textbf{Singular} & \textbf{Plural} \\
\midrule
Nominative & trēo & trēowu \\
Genitive & trēowes & trēowa \\
Dative & trēowe & trēowum \\
Accusative & trēo & trēowu \\
\bottomrule
\end{tabular}
\par\endgroup
% course-panel:T07-P01:end
% mdtable:T07:end

% teaching: CH01-PAR-NEUT-DISYLLABLE
% mdparagraph:P018:begin
Do not apply the monosyllabic \textbf{-u/\allowbreak{}zero} rule mechanically to nouns of two syllables. Learn the following complete patterns. \textbf{Tācen} loses \textbf{e} before endings; \textbf{ǣfen, wæter, nīeten} retain it in the selected forms. The shorter \textbf{wætru} and related forms also occur.[5]
% mdparagraph:P018:end

% mdtable:T08:begin
% course-table-style: 2026-10
% course-panel:T08-P01:begin
\par\begingroup
\fontsize{10.5}{13.2}\selectfont
\setlength{\LTleft}{0pt}\setlength{\LTright}{\fill}
\noindent
\begin{tabular}{@{}>{\raggedright\arraybackslash}p{22.0mm}>{\raggedright\arraybackslash}p{22.688mm}>{\raggedright\arraybackslash}p{22.688mm}>{\raggedright\arraybackslash}p{22.688mm}>{\raggedright\arraybackslash}p{22.688mm}@{}}
\toprule
\textbf{Case} & \textbf{ǣfen singular} & \textbf{ǣfen plural} & \textbf{tācen singular} & \textbf{tācen plural} \\
\midrule
Nominative & ǣfen & ǣfenu & tācen & tācnu \\
Genitive & ǣfenes & ǣfena & tācnes & tācna \\
Dative & ǣfene & ǣfenum & tācne & tācnum \\
Accusative & ǣfen & ǣfenu & tācen & tācnu \\
\bottomrule
\end{tabular}
\par\endgroup
% course-panel:T08-P01:end
% mdtable:T08:end

% mdtable:T09:begin
% course-table-style: 2026-10
% course-panel:T09-P01:begin
\par\begingroup
\fontsize{10.5}{13.2}\selectfont
\setlength{\LTleft}{0pt}\setlength{\LTright}{\fill}
\noindent
\begin{tabular}{@{}>{\raggedright\arraybackslash}p{22.0mm}>{\raggedright\arraybackslash}p{22.688mm}>{\raggedright\arraybackslash}p{22.688mm}>{\raggedright\arraybackslash}p{22.688mm}>{\raggedright\arraybackslash}p{22.688mm}@{}}
\toprule
\textbf{Case} & \textbf{wæter singular} & \textbf{wæter plural} & \textbf{nīeten singular} & \textbf{nīeten plural} \\
\midrule
Nominative & wæter & wæteru & nīeten & nīetenu \\
Genitive & wæteres & wætera & nīetenes & nīetena \\
Dative & wætere & wæterum & nīetene & nīetenum \\
Accusative & wæter & wæteru & nīeten & nīetenu \\
\bottomrule
\end{tabular}
\par\endgroup
% course-panel:T09-P01:end
% mdtable:T09:end

\reserve{4\baselineskip}
% mdparagraph:P019:begin
\section*{5. Strong feminine nouns}\BookHeading{nav-ch01-learner-0005}{subsection}{5. Strong feminine nouns}{1}\BookLabel{bk-ch01-learner-71357d7f7267}
% mdparagraph:P019:end

% teaching: CH01-PAR-FEM-LONG CH01-PAR-FEM-ACC-ZERO
% mdparagraph:P020:begin
The strong feminines required here have no nominative singular ending. The ordinary genitive, dative, and accusative singular end in \textbf{-e}. \textbf{Wyrt} and \textbf{tīd}, however, have an endingless accusative singular in the selected pattern. The selected N/\allowbreak{}A plural is \textbf{-a}; \textbf{-e} is also found.[6]
% mdparagraph:P020:end

% mdtable:T10:begin
% course-table-style: 2026-10
% course-panel:T10-P01:begin
\par\begingroup
\fontsize{10.5}{13.2}\selectfont
\setlength{\LTleft}{0pt}\setlength{\LTright}{\fill}
\noindent
\begin{tabular}{@{}>{\raggedright\arraybackslash}p{22.0mm}>{\raggedright\arraybackslash}p{22.688mm}>{\raggedright\arraybackslash}p{22.688mm}>{\raggedright\arraybackslash}p{22.688mm}>{\raggedright\arraybackslash}p{22.688mm}@{}}
\toprule
\textbf{Case} & \textbf{gesceaft singular} & \textbf{gesceaft plural} & \textbf{wyrt singular} & \textbf{wyrt plural} \\
\midrule
Nominative & gesceaft & gesceafta & wyrt & wyrta \\
Genitive & gesceafte & gesceafta & wyrte & wyrta \\
Dative & gesceafte & gesceaftum & wyrte & wyrtum \\
Accusative & gesceafte & gesceafta & wyrt & wyrta \\
\bottomrule
\end{tabular}
\par\endgroup
% course-panel:T10-P01:end
% mdtable:T10:end

% mdparagraph:P021:begin
\textbf{tīd}\par\nobreak
% mdparagraph:P021:end

% mdtable:T11:begin
% course-table-style: 2026-10
% course-panel:T11-P01:begin
\par\begingroup
\fontsize{10.5}{13.2}\selectfont
\setlength{\LTleft}{0pt}\setlength{\LTright}{\fill}
\noindent
\begin{tabular}{@{}>{\raggedright\arraybackslash}p{39.459mm}>{\raggedright\arraybackslash}p{39.459mm}>{\raggedright\arraybackslash}p{39.459mm}@{}}
\toprule
\textbf{Case} & \textbf{Singular} & \textbf{Plural} \\
\midrule
Nominative & tīd & tīda \\
Genitive & tīde & tīda \\
Dative & tīde & tīdum \\
Accusative & tīd & tīda \\
\bottomrule
\end{tabular}
\par\endgroup
% course-panel:T11-P01:end
% mdtable:T11:end

% teaching: CH01-PAR-FEM-NES CH01-PAR-FEM-UNG
% mdparagraph:P022:begin
Nouns ending in \textbf{-nes} have \textbf{-ness-} before endings: retain the \textbf{e} and double the \textbf{s}. The same pattern governs \textbf{drȳgnes, gelīcnes, fæstnes, nēowolnes, brādnes}. Nouns in \textbf{-ung}, and the assigned \textbf{līehting}, take the endings shown for \textbf{gegaderung}; \textbf{frætwung} follows it. Cook also records singular oblique \textbf{-a} in this group, especially in the dative. The present tables select \textbf{-e}.[6]
% mdparagraph:P022:end

% mdtable:T12:begin
% course-table-style: 2026-10
% course-panel:T12-P01:begin
\par\begingroup
\fontsize{10.5}{13.2}\selectfont
\setlength{\LTleft}{0pt}\setlength{\LTright}{\fill}
\noindent
\begin{tabular}{@{}>{\raggedright\arraybackslash}p{19.0mm}>{\raggedright\arraybackslash}p{23.0mm}>{\raggedright\arraybackslash}p{23.0mm}>{\raggedright\arraybackslash}p{23.0mm}>{\raggedright\arraybackslash}p{24.752mm}@{}}
\toprule
\textbf{Case} & \textbf{andlīcnes singular} & \textbf{andlīcnes plural} & \textbf{gegaderung singular} & \textbf{gegaderung plural} \\
\midrule
Nominative & andlīcnes & andlīcnessa & gegaderung & gegaderunga \\
Genitive & andlīcnesse & andlīcnessa & gegaderunge & gegaderunga \\
Dative & andlīcnesse & andlīcnessum & gegaderunge & gegaderungum \\
Accusative & andlīcnesse & andlīcnessa & gegaderunge & gegaderunga \\
\bottomrule
\end{tabular}
\par\endgroup
% course-panel:T12-P01:end
% mdtable:T12:end

\reserve{4\baselineskip}
% mdparagraph:P023:begin
\section*{6. Weak nouns}\BookHeading{nav-ch01-learner-0006}{subsection}{6. Weak nouns}{1}\BookLabel{bk-ch01-learner-9c5fd7a11f03}
% mdparagraph:P023:end

% teaching: CH01-PAR-WEAK
% mdparagraph:P024:begin
Weak masculine \textbf{steorra} and weak feminine \textbf{eorðe} differ in the nominative singular. Both have \textbf{-an} in the other singular cases and the N/\allowbreak{}A plural, \textbf{-ena} in the genitive plural, and \textbf{-um} in the dative plural. Replace nominative \textbf{-a/\allowbreak{}-e} with these endings.[7]
% mdparagraph:P024:end

% mdtable:T13:begin
% course-table-style: 2026-10
% course-panel:T13-P01:begin
\par\begingroup
\fontsize{10.5}{13.2}\selectfont
\setlength{\LTleft}{0pt}\setlength{\LTright}{\fill}
\noindent
\begin{tabular}{@{}>{\raggedright\arraybackslash}p{22.0mm}>{\raggedright\arraybackslash}p{22.688mm}>{\raggedright\arraybackslash}p{22.688mm}>{\raggedright\arraybackslash}p{22.688mm}>{\raggedright\arraybackslash}p{22.688mm}@{}}
\toprule
\textbf{Case} & \textbf{steorra singular} & \textbf{steorra plural} & \textbf{eorðe singular} & \textbf{eorðe plural} \\
\midrule
Nominative & steorra & steorran & eorðe & eorðan \\
Genitive & steorran & steorrena & eorðan & eorðena \\
Dative & steorran & steorrum & eorðan & eorðum \\
Accusative & steorran & steorran & eorðan & eorðan \\
\bottomrule
\end{tabular}
\par\endgroup
% course-panel:T13-P01:end
% mdtable:T13:end

% mdparagraph:P025:begin
The plural of \textbf{eorðe} is displayed for recognizing its pattern; the exercises use singular “earth.” An \textbf{-an} form alone cannot identify either case or number.
% mdparagraph:P025:end

\reserve{4\baselineskip}
% mdparagraph:P026:begin
\section*{7. Individual paradigms}\BookHeading{nav-ch01-learner-0007}{subsection}{7. Individual paradigms}{1}\BookLabel{bk-ch01-learner-3d6b15e7b6bf}
% mdparagraph:P026:end

% teaching: CH01-PAR-EXCEPTION-MANN CH01-PAR-EXCEPTION-NIHT
% mdparagraph:P027:begin
\textbf{Mann} has \textbf{męnn} in the dative singular and N/\allowbreak{}A plural. In the selected \textbf{niht} paradigm, every singular case and the N/\allowbreak{}A plural are \textbf{niht}; the genitive and dative plurals are \textbf{nihta, nihtum}. Genitive/\allowbreak{}dative singular \textbf{nihte} is a valid alternative, but is not the form selected for recall here.[8]
% mdparagraph:P027:end

% mdtable:T14:begin
% course-table-style: 2026-10
% course-panel:T14-P01:begin
\par\begingroup
\fontsize{10.5}{13.2}\selectfont
\setlength{\LTleft}{0pt}\setlength{\LTright}{\fill}
\noindent
\begin{tabular}{@{}>{\raggedright\arraybackslash}p{22.0mm}>{\raggedright\arraybackslash}p{22.688mm}>{\raggedright\arraybackslash}p{22.688mm}>{\raggedright\arraybackslash}p{22.688mm}>{\raggedright\arraybackslash}p{22.688mm}@{}}
\toprule
\textbf{Case} & \textbf{mann singular} & \textbf{mann plural} & \textbf{niht singular} & \textbf{niht plural} \\
\midrule
Nominative & mann & męnn & niht & niht \\
Genitive & mannes & manna & niht & nihta \\
Dative & męnn & mannum & niht & nihtum \\
Accusative & mann & męnn & niht & niht \\
\bottomrule
\end{tabular}
\par\endgroup
% course-panel:T14-P01:end
% mdtable:T14:end

% teaching: CH01-PAR-EXCEPTION-SAE CH01-REF-SAE
% mdparagraph:P028:begin
\textbf{Sǣ} is masculine in the selected vocabulary. Its nominative plural is \textbf{sǣs}, not a form in \textbf{-as}. The starred genitive plural below is a reconstruction: consult it only as a reference, and do not use it in these exercises.[8]
% mdparagraph:P028:end

% mdparagraph:P029:begin
\textbf{sǣ}\par\nobreak
% mdparagraph:P029:end

% mdtable:T15:begin
% course-table-style: 2026-10
% course-panel:T15-P01:begin
\par\begingroup
\fontsize{10.5}{13.2}\selectfont
\setlength{\LTleft}{0pt}\setlength{\LTright}{\fill}
\noindent
\begin{tabular}{@{}>{\raggedright\arraybackslash}p{39.459mm}>{\raggedright\arraybackslash}p{39.459mm}>{\raggedright\arraybackslash}p{39.459mm}@{}}
\toprule
\textbf{Case} & \textbf{Singular} & \textbf{Plural} \\
\midrule
Nominative & sǣ & sǣs \\
Genitive & sǣs & *sǣwa — reconstructed; not tested \\
Dative & sǣ & sǣum \\
Accusative & sǣ & sǣs \\
\bottomrule
\end{tabular}
\par\endgroup
% course-panel:T15-P01:end
% mdtable:T15:end

% teaching: CH01-PAR-EXCEPTION-THIESTRU
% mdparagraph:P030:begin
\textbf{Ðīestru}, “darkness,” is plural in Old English. Its complete assigned paradigm is plural only.
% mdparagraph:P030:end

% mdparagraph:P031:begin
\textbf{ðīestru}\par\nobreak
% mdparagraph:P031:end

% mdtable:T16:begin
% course-table-style: 2026-10
% course-panel:T16-P01:begin
\par\begingroup
\fontsize{10.5}{13.2}\selectfont
\setlength{\LTleft}{0pt}\setlength{\LTright}{\fill}
\noindent
\begin{tabular}{@{}>{\raggedright\arraybackslash}p{60.594mm}>{\raggedright\arraybackslash}p{60.594mm}@{}}
\toprule
\textbf{Case} & \textbf{Plural} \\
\midrule
Nominative & ðīestru \\
Genitive & ðīestra \\
Dative & ðīestrum \\
Accusative & ðīestru \\
\bottomrule
\end{tabular}
\par\endgroup
% course-panel:T16-P01:end
% mdtable:T16:end

\reserve{4\baselineskip}
% mdparagraph:P032:begin
\section*{8. Subject pronouns and the indicative copula}\BookHeading{nav-ch01-learner-0008}{subsection}{8. Subject pronouns and the indicative copula}{1}
% mdparagraph:P032:end

% teaching: CH01-GRAM-SUBJECT-PRONOUNS CH01-PAR-SUBJECTS CH01-VOC-SUBJECT-ic CH01-VOC-SUBJECT-thu CH01-VOC-SUBJECT-he
% mdparagraph:P033:begin
The first person is the speaker; the second is the person addressed; the third is the person or thing spoken about. Learn these complete \textbf{nominative subject forms} with the verb “be.” The other cases of the personal pronouns follow in Unit 2.[16]
% mdparagraph:P033:end

% mdtable:T17:begin
% course-table-style: 2026-10
% course-panel:T17-P01:begin
\par\begingroup
\fontsize{10.5}{13.2}\selectfont
\setlength{\LTleft}{0pt}\setlength{\LTright}{\fill}
\noindent
\begin{tabular}{@{}>{\raggedright\arraybackslash}p{21.929mm}>{\raggedright\arraybackslash}p{37.279mm}>{\raggedright\arraybackslash}p{25.218mm}>{\raggedright\arraybackslash}p{31.139mm}@{}}
\toprule
\textbf{Person}\BookPoint{bk-ch01-learner-f4792afbd8d6} & \textbf{Singular} & \textbf{Dual: exactly two} & \textbf{Plural} \\
\midrule
First & ic — I & wit — we two & wē — we \\
Second & ðū — you & git — you two & gē — you \\
Third & hē — he; hit — it; hēo — she & — & hīe — they, all genders \\
\bottomrule
\end{tabular}
\par\endgroup
% course-panel:T17-P01:end
% mdtable:T17:end

% teaching: CH01-GRAM-SUBJECT-REFERENCE
% mdparagraph:P034:begin
In\BookPoint{bk-ch01-learner-8f8ad9d5a04c} the third-person singular, \textbf{hē, hit, hēo} normally refer to masculine, neuter, and feminine nouns respectively. Use grammatical gender when referring to a thing: a feminine noun need not name a female. All plural nouns take \textbf{hīe}. The context identifies which noun the pronoun refers to; gender and number help to check that reference.
% mdparagraph:P034:end

% teaching: CH01-GRAM-COPULA
% mdparagraph:P035:begin
These are the complete present \textbf{eom} and preterite \textbf{wæs} series of the indicative “be.” The preterite expresses a past state. Other present forms, negative forms, moods, and nonfinite constructions will be treated separately.[9]
% mdparagraph:P035:end

% mdtable:T18:begin
% course-table-style: 2026-10
% course-panel:T18-P01:begin
\par\begingroup
\fontsize{10.5}{13.2}\selectfont
\setlength{\LTleft}{0pt}\setlength{\LTright}{\fill}
\noindent
\begin{tabular}{@{}>{\raggedright\arraybackslash}p{39.459mm}>{\raggedright\arraybackslash}p{39.459mm}>{\raggedright\arraybackslash}p{39.459mm}@{}}
\toprule
\textbf{Person and number}\BookPoint{bk-ch01-learner-814e0a394df4} & \textbf{Present} & \textbf{Preterite} \\
\midrule
1 sg.: I & eom & wæs \\
2 sg.: you & eart & wǣre \\
3 sg.: he, she, it & is & wæs \\
1–3 pl.: we, you, they & sindon & wǣron \\
\bottomrule
\end{tabular}
\par\endgroup
% course-panel:T18-P01:end
% mdtable:T18:end

% teaching: CH01-GRAM-DUAL-SUBJECTS
% mdparagraph:P036:begin
\textbf{1–3 pl.}\BookPoint{bk-ch01-learner-f73641fc4bb6} gives the shared form for \textbf{wē, gē, hīe}. The dual subjects \textbf{wit} and \textbf{git} also take these ordinary plural verb forms: there is no separate dual form of the copula. Use a dual when specifically saying “we two” or “you two”; the ordinary plurals do not exclude a group of two. Cook also gives \textbf{sind, sint} beside the selected \textbf{sindon}.[16]
% mdparagraph:P036:end

\reserve{4\baselineskip}
% mdparagraph:P037:begin
\section*{9. Basic constructions}\BookHeading{nav-ch01-learner-0009}{subsection}{9. Basic constructions}{1}\BookLabel{bk-ch01-learner-086b075156aa}
% mdparagraph:P037:end

% teaching: CH01-GRAM-SUBJECT-PREDICATE CH01-GRAM-AGREEMENT
% mdparagraph:P038:begin
The subject is nominative. A noun that identifies it after “be” is also nominative; English \textbf{a/\allowbreak{}an} need not receive a separate word in such a predicate. The copula agrees with the subject in person and number.[10]
% mdparagraph:P038:end

% example: CH01-EXAMPLE07
% mdparagraph:P039:begin
\textbf{Ðā hwalas sindon dēor.} — The whales are animals.
% mdparagraph:P039:end

% mdparagraph:P040:begin
Both subject \textbf{hwalas} and predicate \textbf{dēor} are nominative plural; the neuter plural \textbf{dēor} has no added ending.
% mdparagraph:P040:end

% example: CH01-EXAMPLE09
% mdparagraph:P041:begin
\textbf{Ic eom mann.} — I am a human being.
% mdparagraph:P041:end

% example: CH01-EXAMPLE10
% mdparagraph:P042:begin
\textbf{Wit sindon męnn.} — We two are human beings.
% mdparagraph:P042:end

% teaching: CH01-GRAM-COORDINATION
% mdparagraph:P043:begin
\textbf{And} joins words or phrases. When definite noun phrases are joined, each article agrees with its own noun. Gender belongs to that noun; it is not taken from the other noun in the pair.
% mdparagraph:P043:end

% teaching: CH01-GRAM-ON
% mdparagraph:P044:begin
For the location of an entity, use \textbf{on + dative}, “in/\allowbreak{}on.”
% mdparagraph:P044:end

% example: CH01-EXAMPLE01
% mdparagraph:P045:begin
\textbf{Se fisc is on ðǣm wætere.} — The fish is in the water.
% mdparagraph:P045:end

% example: CH01-EXAMPLE02
% mdparagraph:P046:begin
\textbf{Ðā męnn wǣron on ðǣm hofum.} — The men were in the dwellings.
% mdparagraph:P046:end

% teaching: CH01-GRAM-GENITIVE
% mdparagraph:P047:begin
A dependent genitive expresses possession or a relation such as the source, contents, or kind of something. In the phrase pattern used for Cook's Exercise 3, place the genitive first and omit the article of the \textbf{head noun}, while retaining the article belonging to the dependent. The head remains in the case required by its own role. A proper name such as \textbf{God} takes no article here.[11]
% mdparagraph:P047:end

% example: CH01-EXAMPLE03
% mdparagraph:P048:begin
\textbf{ðæs mannes fugol} — the man's bird
% mdparagraph:P048:end

% mdparagraph:P049:begin
Here \textbf{ðæs mannes} is the dependent genitive; \textbf{fugol} is the head, nominative in the isolated naming phrase. The construction does not mean that every English “the” disappears.
% mdparagraph:P049:end

% teaching: CH01-GRAM-NESTED-GENITIVE
% mdparagraph:P050:begin
The\BookPoint{bk-ch01-learner-0b109d29614d} same dependent may stand inside a prepositional phrase. Its genitive does not change the case required for the head.
% mdparagraph:P050:end

% example: CH01-EXAMPLE08
% mdparagraph:P051:begin
\textbf{on ðæs fugles trēowe} — in the bird's tree
% mdparagraph:P051:end

% mdparagraph:P052:begin
\textbf{Trēowe} is dative after locative \textbf{on}; \textbf{ðæs fugles} remains genitive. \textbf{Ðæs} agrees with \textbf{fugles}, not with \textbf{trēowe}.
% mdparagraph:P052:end

% teaching: CH01-GRAM-OBJECT
% mdparagraph:P053:begin
The\BookPoint{bk-ch01-learner-422453d6be7d} direct object is accusative. In the following example, \textbf{geseah} is supplied: “saw,” third-person singular preterite. Its conjugation is not required here.[12]
% mdparagraph:P053:end

% example: CH01-EXAMPLE04
% mdparagraph:P054:begin
\textbf{Se mann geseah ðone fisc.} — The man saw the fish.
% mdparagraph:P054:end

% mdparagraph:P055:begin
\textbf{Mann} and \textbf{fisc} have identical nominative and accusative singular spellings, but their articles distinguish subject \textbf{se} from object \textbf{ðone}.
% mdparagraph:P055:end

% teaching: CH01-GRAM-RECIPIENT CH01-FRAME-GIVING
% mdparagraph:P056:begin
The\BookPoint{bk-ch01-learner-d64391595e92} recipient of a gift is dative; the thing given is accusative. The supplied \textbf{sealde} means “gave,” third-person singular preterite.[12]
% mdparagraph:P056:end

% example: CH01-EXAMPLE05
% mdparagraph:P057:begin
\textbf{God sealde ðǣm fugle męte.} — God gave the bird food.
% mdparagraph:P057:end

% mdparagraph:P058:begin
\textbf{Ðǣm fugle} is the recipient; \textbf{męte} is the direct object. English may express the recipient with “to,” but this Old English construction uses the dative without a preposition.
% mdparagraph:P058:end

% teaching: CH01-GRAM-FRAM
% mdparagraph:P059:begin
Use\BookPoint{bk-ch01-learner-2b2114170c4e} \textbf{fram + dative} for what something is separated from. The supplied \textbf{gedǣlde} means “separated/\allowbreak{}divided,” third-person singular preterite.[12]
% mdparagraph:P059:end

% example: CH01-EXAMPLE06
% mdparagraph:P060:begin
\textbf{Se mann gedǣlde ðæt sǣd fram ðǣm gærse.} — The man separated the seed from the grass.
% mdparagraph:P060:end

% mdparagraph:P061:begin
The thing separated, \textbf{ðæt sǣd}, is accusative; \textbf{ðǣm gærse} is dative after \textbf{fram}. These are particular constructions. They do not establish one case rule for all spatial prepositions.
% mdparagraph:P061:end

\reserve{4\baselineskip}
% mdparagraph:P062:begin
\section*{10. Vocabulary}\BookHeading{nav-ch01-learner-0010}{subsection}{10. Vocabulary}{1}
% mdparagraph:P062:end

% mdparagraph:P063:begin
Each pattern reference gives a section and its model noun: \textbf{§4: hof}, for example, means the \textbf{hof} pattern in §4. References to \textbf{sǣd}, \textbf{trēo}, \textbf{tācen}, or another named neuter select that particular pattern. They are not interchangeable.
% mdparagraph:P063:end

% mdparagraph:P064:begin
Learn each noun with its gender and diagnostic forms. \textbf{G sg} identifies the singular stem; \textbf{N pl} identifies the plural pattern. The full tables and rules above supply the other cases. \textbf{S} marks a noun used only in the singular in composition here: its plural morphology is for recognition, with no additional plural sense assigned. \textbf{P} marks plural-only \textbf{ðīestru}. The singular divine name \textbf{God} has no plural assigned. Abbreviations: \textbf{m.}, masculine; \textbf{f.}, feminine; \textbf{n.}, neuter.[13]
% mdparagraph:P064:end

% mdtable:T19:begin
% course-table-style: 2026-10
% teaching: CH01-VOC-aefen
% teaching: CH01-VOC-morgen
% teaching: CH01-VOC-daeg
% teaching: CH01-VOC-niht
% teaching: CH01-VOC-eorthe
% teaching: CH01-VOC-heofon
% teaching: CH01-VOC-fisc
% teaching: CH01-VOC-fugol
% teaching: CH01-VOC-gaers
% teaching: CH01-VOC-saed
% teaching: CH01-VOC-treo
% teaching: CH01-VOC-wyrt
% teaching: CH01-VOC-drygnes
% teaching: CH01-VOC-waeter
% teaching: CH01-VOC-andlicnes
% teaching: CH01-VOC-gelicnes
% teaching: CH01-VOC-leoht
% teaching: CH01-VOC-thiestru
% teaching: CH01-VOC-tacen
% teaching: CH01-VOC-tid
% teaching: CH01-VOC-gast
% teaching: CH01-VOC-lif
% teaching: CH01-VOC-gesceaft
% teaching: CH01-VOC-mete
% teaching: CH01-VOC-gear
% teaching: CH01-VOC-mann
% teaching: CH01-VOC-deor
% teaching: CH01-VOC-waestm
% teaching: CH01-VOC-faestnes
% teaching: CH01-VOC-neowolnes
% teaching: CH01-VOC-sae
% teaching: CH01-VOC-hwael
% teaching: CH01-VOC-nieten
% teaching: CH01-VOC-anginn
% teaching: CH01-VOC-god
% teaching: CH01-VOC-bradnes
% teaching: CH01-VOC-gegaderung
% teaching: CH01-VOC-lyft
% teaching: CH01-VOC-fraetwung
% teaching: CH01-VOC-weorc
% teaching: CH01-VOC-cynn
% teaching: CH01-VOC-steorra
% teaching: CH01-VOC-liehting
% teaching: CH01-VOC-hof
% course-panel:T19-P01:begin
\par\begingroup
\fontsize{10.5}{13.2}\selectfont
\setlength{\LTleft}{0pt}\setlength{\LTright}{\fill}
\begin{longtable}{@{}>{\raggedright\arraybackslash}p{24.0mm}>{\raggedright\arraybackslash}p{39.0mm}>{\raggedright\arraybackslash}p{28.0mm}>{\raggedright\arraybackslash}p{24.564mm}@{}}
\toprule
\textbf{Noun; gender} & \textbf{Meaning} & \textbf{Diagnostic forms} & \textbf{Pattern; limit} \\
\midrule
\endfirsthead
\toprule
\textbf{Noun; gender} & \textbf{Meaning} & \textbf{Diagnostic forms} & \textbf{Pattern; limit} \\
\midrule
\endhead
\midrule
\multicolumn{4}{r}{\footnotesize Continued on the following page.}\\
\endfoot
\bottomrule
\endlastfoot
 ǣfen (n.) & evening & G sg ǣfenes N pl ǣfenu & §4: ǣfen \\
 morgen (m.) & morning & G sg morgnes N pl morgnas & §3: morgen \\
 dæg (m.) & day & G sg dæges N pl dagas & §3: dæg \\
 niht (f.) & night & G sg niht N pl niht & §7: niht \\
 eorðe (f.) & earth & G sg eorðan N pl eorðan & §6: eorðe; S \\
 heofon (m.) & heaven; sky & G sg heofones N pl heofonas & §3: fisc \\
 fisc (m.) & fish & G sg fisces N pl fiscas & §3: fisc \\
 fugol (m.) & bird & G sg fugles N pl fuglas & §3: fugol \\
 gærs (n.) & grass & G sg gærses N pl gærs & §4: sǣd; S \\
 sǣd (n.) & seed & G sg sǣdes N pl sǣd & §4: sǣd \\
 trēo (n.) & tree & G sg trēowes N pl trēowu & §4: trēo \\
 wyrt (f.) & herb; plant & G sg wyrte N pl wyrta & §5: wyrt \\
 drȳgnes (f.) & dry land & G sg drȳgnesse N pl drȳgnessa & §5: andlīcnes; S \\
 wæter (n.) & water & G sg wæteres N pl wæteru & §4: wæter \\
 andlīcnes (f.) & image & G sg andlīcnesse N pl andlīcnessa & §5: andlīcnes \\
 gelīcnes (f.) & likeness & G sg gelīcnesse N pl gelīcnessa & §5: andlīcnes \\
 lēoht (n.) & light & G sg lēohtes N pl lēoht & §4: sǣd; S \\
 ðīestru (P) & darkness & G sg — N pl ðīestru & §7: ðīestru; P; G pl ðīestra \\
 tācen (n.) & sign & G sg tācnes N pl tācnu & §4: tācen \\
 tīd (f.) & season; time & G sg tīde N pl tīda & §5: tīd \\
 gāst (m.) & spirit & G sg gāstes N pl gāstas & §3: fisc \\
 līf (n.) & life & G sg līfes N pl līf & §4: sǣd; S \\
 gesceaft (f.) & creature; created being or thing & G sg gesceafte N pl gesceafta & §5: gesceaft \\
 męte (m.) & food & G sg mętes N pl mętas & §3: męte; S \\
 gēar (n.) & year & G sg gēares N pl gēar & §4: sǣd \\
 mann (m.) & man; human being & G sg mannes N pl męnn & §7: mann \\
 dēor (n.) & animal; beast & G sg dēores N pl dēor & §4: sǣd \\
 wæstm (m.) & fruit; produce & G sg wæstmes N pl wæstmas & §3: fisc \\
 fæstnes (f.) & firmament & G sg fæstnesse N pl fæstnessa & §5: andlīcnes; S \\
 nēowolnes (f.) & abyss; the deep & G sg nēowolnesse N pl nēowolnessa & §5: andlīcnes; S \\
 sǣ (m.) & sea & G sg sǣs N pl sǣs & §7: sǣ \\
 hwæl (m.) & whale & G sg hwæles N pl hwalas & §3: hwæl \\
 nīeten (n.) & animal; beast; in the plural, cattle & G sg nīetenes N pl nīetenu & §4: nīeten \\
 anginn (n.) & beginning & G sg anginnes N pl anginn & §4: sǣd \\
 God (m.) & God (proper divine name) & G sg Godes N pl — & §3: fisc; divine name, singular \\
 brādnes (f.) & surface; face (of the deep) & G sg brādnesse N pl brādnessa & §5: andlīcnes; S \\
 gegaderung (f.) & gathering; collection; congregation (of waters) & G sg gegaderunge N pl gegaderunga & §5: gegaderung \\
 lyft (n.) & air & G sg lyftes N pl lyft & §4: sǣd; S \\
 frætwung (f.) & array & G sg frætwunge N pl frætwunga & §5: gegaderung; S \\
 weorc (n.) & work & G sg weorces N pl weorc & §4: sǣd \\
 cynn (n.) & kind & G sg cynnes N pl cynn & §4: sǣd \\
 steorra (m.) & star & G sg steorran N pl steorran & §6: steorra \\
 līehting (f.) & illumination & G sg līehtinge N pl līehtinga & §5: gegaderung; S \\
 hof (n.) & dwelling & G sg hofes N pl hofu & §4: hof \\
\end{longtable}
\par\endgroup
% course-panel:T19-P01:end
% mdtable:T19:end

% mdparagraph:P065:begin
Pronunciation: \textbf{cynn} begins with \textbf{k} as in \emph{kit}; \textbf{gærs} begins with stopped \textbf{g} as in \emph{go}, despite its following \textbf{æ}. Learn these supplied sounds with the words.[17]
% mdparagraph:P065:end

% mdparagraph:P066:begin
The other words required for composition are \textbf{and}, “and”; \textbf{on + dative}, “in/\allowbreak{}on” for location; \textbf{fram + dative}, “from” in separation; the article; and the subject pronouns and indicative forms in §8. Verbs printed in frames and the additional help for the reading are supplied, not additional vocabulary to recall independently.
% mdparagraph:P066:end

\reserve{4\baselineskip}
% mdparagraph:P067:begin
\section*{11. Forms and agreement}\BookHeading{nav-ch01-learner-0011}{subsection}{11. Forms and agreement}{1}
% mdparagraph:P067:end

% mdparagraph:P068:begin
The tasks in §§11–12 and §§14–17 are new practice. Use the selected patterns; the key distinguishes valid source variants from the particular forms requested. In the noun tasks, translate only the marked words, not the rest of the English context. Choose the case from the word's function.
% mdparagraph:P068:end

% item: CH01-F01
% mdparagraph:P069:begin
\itempara{1}{Write the bold noun phrase in Old English and label its case: (a) \textbf{The fish} is food; (b) the body of \textbf{the fish}; (c) God gave food to \textbf{the fish}; (d) the bird saw \textbf{the fish}. First make each fish singular, then rewrite all four phrases for several fish. Include the article.}
% mdparagraph:P069:end

% item: CH01-F02
% mdparagraph:P070:begin
\itempara{2}{Write only the bold noun, without an article: (a) several \textbf{dwellings} are on the earth; (b) the man saw several \textbf{dwellings}; (c) the \textbf{seeds} are food; (d) the bird saw the \textbf{seeds}. Identify each case and number. Explain why the two nouns differ in their plural endings although their roles are parallel.}
% mdparagraph:P070:end

% item: CH01-F03
% mdparagraph:P071:begin
\itempara{3}{Give only the bold noun, without an article, and explain its case: (a) the length of several \textbf{days}; (b) the food of several \textbf{whales}; (c) God gave food to one \textbf{man}; (d) the darkness of several \textbf{nights}. Preserve each noun's own stem changes.}
% mdparagraph:P071:end

% item: CH01-F04
% mdparagraph:P072:begin
\itempara{4}{Give only the bold noun, without an article: (a) the fruit of one \textbf{tree}; (b) the bird is in one \textbf{tree}; (c) several \textbf{trees} are in the dwelling; (d) the meaning of one \textbf{sign}; (e) several \textbf{signs} are in the sky; (f) the sound of the \textbf{water}; (g) the \textbf{waters} are in the sea; (h) several \textbf{evenings} have passed; (i) the \textbf{cattle} are in the dwelling. Label the forms. Explain the \textbf{w} in (c), and contrast the loss or retention of the middle stem vowel in (e) and (g).}
% mdparagraph:P072:end

% item: CH01-F05
% mdparagraph:P073:begin
\itempara{5}{Give only the bold noun, without an article: (a) \textbf{the star} is in the sky; (b) the light of \textbf{the star}; (c) the man is looking at \textbf{the star}, as a direct object; (d) the bird is under \textbf{the star}, using a supplied dative frame. First use one star, then several stars. Next give singular \textbf{earth} as (e) subject, (f) possessor in “the surface of the earth,” (g) location after \textbf{on}, and (h) direct object in “God created the earth.” Explain why an isolated \textbf{-an} form cannot settle the analysis. Only the noun forms are requested; the verbs and the English prepositions in (c–d) are not to be translated.}
% mdparagraph:P073:end

% item: CH01-F06
% mdparagraph:P074:begin
\itempara{6}{Give only the bold noun, without an article: (a) the man saw one \textbf{herb}; (b) the seed of one \textbf{herb}; (c) the man marked one \textbf{season} as the direct object; (d) the beginning of one \textbf{season}; (e) the man saw one \textbf{creature}. Explain why the direct objects do not all take the same feminine ending.}
% mdparagraph:P074:end

% item: CH01-F07
% mdparagraph:P075:begin
\itempara{7}{Give only the bold noun, without an article: (a) in the \textbf{firmament}, using locative \textbf{on}; (b) the form of several \textbf{images}; (c) several \textbf{likenesses} are in the dwelling; (d) in several \textbf{gatherings}, using locative \textbf{on}. Identify the case and number, and explain the stem before each ending.}
% mdparagraph:P075:end

% item: CH01-F08
% mdparagraph:P076:begin
\itempara{8}{Reference check: locate \textbf{ðȳ} and \textbf{ðon} in the article table. Identify their case and singular genders. This checks recognition of the reference forms; their use in composition is not required.}
% mdparagraph:P076:end

\reserve{4\baselineskip}
% mdparagraph:P077:begin
\subsection*{The copula}\BookHeading{nav-ch01-learner-0012}{subsubsection}{The copula}{0}
% mdparagraph:P077:end

% mdparagraph:P078:begin
Write complete Old English sentences using a subject pronoun. Preserve the intended person, number, time, and location. Use a dual wherever the prompt explicitly says “two.”
% mdparagraph:P078:end

% item: CH01-COP01
% mdparagraph:P079:begin
\itempara{1}{You are speaking about yourself: “I was in the tree.” Then change only the time: “I am in the tree.”}
% mdparagraph:P079:end

% item: CH01-COP02
% mdparagraph:P080:begin
\itempara{2}{Address one person: “You are in the water.” Then place the same state in the past.}
% mdparagraph:P080:end

% item: CH01-COP03
% mdparagraph:P081:begin
\itempara{3}{Replace the subject of each sentence with its appropriate pronoun and translate: (a) “The spirit was in the air”; (b) “The sign is in the dwelling”; (c) “The image is on the earth.” Which choice depends on grammatical gender rather than sex?}
% mdparagraph:P081:end

% item: CH01-COP04
% mdparagraph:P082:begin
\itempara{4}{You speak for yourself and one companion: “We two are in the sea.” Put this in the past, then restate the past sentence for a group of three that includes you. Explain which pronoun changes and why the verb's number does not.}
% mdparagraph:P082:end

% item: CH01-COP05
% mdparagraph:P083:begin
\itempara{5}{Address two companions: “You two were in the dwelling.” Put this in the present, then restate the present sentence to address a group of three. Explain the choices of pronoun and verb.}
% mdparagraph:P083:end

% item: CH01-COP06
% mdparagraph:P084:begin
\itempara{6}{The whales have already been mentioned. Say “They were in the water,” then put this in the present. Would the subject pronoun change if it referred instead to several feminine \textbf{gesceafta}, “creatures”? Explain.}
% mdparagraph:P084:end

\reserve{4\baselineskip}
% mdparagraph:P085:begin
\section*{12. Case in sentences}\BookHeading{nav-ch01-learner-0013}{subsection}{12. Case in sentences}{1}
% mdparagraph:P085:end

% mdparagraph:P086:begin
Complete each frame with the noun and article phrases needed for the English sentence. The printed verb is supplied in its exact form; choose the noun cases, numbers, and articles. The prepositions in the frames use the constructions of §9.
% mdparagraph:P086:end

\reserve{7\baselineskip}
% item: CH01-CASE01
% mdparagraph:P087:begin
\itempara{1}{The man saw the bird.}
% mdparagraph:P087:end

% mdparagraph:P088:begin
\textbf{\_\_\_\_\_ geseah \_\_\_\_\_.}\par\nobreak
% mdparagraph:P088:end

% mdparagraph:P089:begin
Supplied: \textbf{geseah}, “saw,” third-person singular preterite.
% mdparagraph:P089:end

\reserve{7\baselineskip}
% item: CH01-CASE02
% mdparagraph:P090:begin
\itempara{2}{The bird saw the man.}
% mdparagraph:P090:end

% mdparagraph:P091:begin
\textbf{\_\_\_\_\_ geseah \_\_\_\_\_.}\par\nobreak
% mdparagraph:P091:end

% mdparagraph:P092:begin
Supplied: \textbf{geseah}, “saw,” third-person singular preterite.
% mdparagraph:P092:end

\reserve{7\baselineskip}
% item: CH01-CASE03
% mdparagraph:P093:begin
\itempara{3}{God created the heavens and the earth. Use masculine heofon for 'heavens'.}
% mdparagraph:P093:end

% mdparagraph:P094:begin
\textbf{\_\_\_\_\_ gescēop \_\_\_\_\_ and \_\_\_\_\_.}\par\nobreak
% mdparagraph:P094:end

% mdparagraph:P095:begin
Supplied: \textbf{gescēop}, “created,” third-person singular preterite.
% mdparagraph:P095:end

\reserve{7\baselineskip}
% item: CH01-CASE04
% mdparagraph:P096:begin
\itempara{4}{The man saw the stars.}
% mdparagraph:P096:end

% mdparagraph:P097:begin
\textbf{\_\_\_\_\_ geseah \_\_\_\_\_.}\par\nobreak
% mdparagraph:P097:end

% mdparagraph:P098:begin
Supplied: \textbf{geseah}, “saw,” third-person singular preterite.
% mdparagraph:P098:end

\reserve{7\baselineskip}
% item: CH01-CASE05
% mdparagraph:P099:begin
\itempara{5}{The man saw the fish in the sea.}
% mdparagraph:P099:end

% mdparagraph:P100:begin
\textbf{\_\_\_\_\_ geseah \_\_\_\_\_ on \_\_\_\_\_.}\par\nobreak
% mdparagraph:P100:end

% mdparagraph:P101:begin
Supplied: \textbf{geseah}, “saw,” third-person singular preterite. The fish is in the sea.
% mdparagraph:P101:end

\reserve{7\baselineskip}
% item: CH01-CASE06
% mdparagraph:P102:begin
\itempara{6}{God separated the light from the darkness.}
% mdparagraph:P102:end

% mdparagraph:P103:begin
\textbf{\_\_\_\_\_ gedǣlde \_\_\_\_\_ fram \_\_\_\_\_.}\par\nobreak
% mdparagraph:P103:end

% mdparagraph:P104:begin
Supplied: \textbf{gedǣlde}, “separated/\allowbreak{}divided,” third-person singular preterite.
% mdparagraph:P104:end

\reserve{7\baselineskip}
% item: CH01-CASE07
% mdparagraph:P105:begin
\itempara{7}{The bird saw the seed in the man's dwelling.}
% mdparagraph:P105:end

% mdparagraph:P106:begin
\textbf{\_\_\_\_\_ geseah \_\_\_\_\_ on \_\_\_\_\_.}\par\nobreak
% mdparagraph:P106:end

% mdparagraph:P107:begin
Supplied: \textbf{geseah}, “saw,” third-person singular preterite. The seed is in the dwelling.
% mdparagraph:P107:end

\reserve{7\baselineskip}
% item: CH01-CASE08
% mdparagraph:P108:begin
\itempara{8}{The bird saw the man in the dwelling.}
% mdparagraph:P108:end

% mdparagraph:P109:begin
\textbf{\_\_\_\_\_ geseah \_\_\_\_\_ on \_\_\_\_\_.}\par\nobreak
% mdparagraph:P109:end

% mdparagraph:P110:begin
Supplied: \textbf{geseah}, “saw,” third-person singular preterite. The man is in the dwelling.
% mdparagraph:P110:end

\reserve{9\baselineskip}
% item: CH01-CASE09
% mdparagraph:P111:begin
\itempara{9}{Complete the printed sealde frame to express: God gave food to the man. Identify the recipient and the thing given.}
% mdparagraph:P111:end

% mdparagraph:P112:begin
\textbf{\_\_\_\_\_ sealde \_\_\_\_\_ \_\_\_\_\_.}\par\nobreak
% mdparagraph:P112:end

% mdparagraph:P113:begin
Supplied: \textbf{sealde}, “gave,” third-person singular preterite. Put the recipient before the thing given in this frame.
% mdparagraph:P113:end

\reserve{9\baselineskip}
% item: CH01-CASE10
% mdparagraph:P114:begin
\itempara{10}{Complete the printed sealde frame to express: The man gave seed to the birds. Identify the recipient and the thing given.}
% mdparagraph:P114:end

% mdparagraph:P115:begin
\textbf{\_\_\_\_\_ sealde \_\_\_\_\_ \_\_\_\_\_.}\par\nobreak
% mdparagraph:P115:end

% mdparagraph:P116:begin
Supplied: \textbf{sealde}, “gave,” third-person singular preterite. Put the recipient before the thing given in this frame.
% mdparagraph:P116:end

\reserve{4\baselineskip}
% mdparagraph:P117:begin
\section*{13. Cook's noun-phrase exercises}\BookHeading{nav-ch01-learner-0014}{subsection}{13. Cook's noun-phrase exercises}{1}
% mdparagraph:P117:end

% mdparagraph:P118:begin
Translate the following original prompts using the vocabulary of this unit. In Exercises 1–2, make the isolated phrases nominative and preserve the requested number. In Exercise 3, keep the head nominative; place its genitive dependent first, omitting the head's own article as explained in §9. The dependent retains its appropriate article. Use \textbf{dēor} for “animal/\allowbreak{}beast” and \textbf{nīeten} for plural “cattle”; \textbf{brādnes} means the surface in “face of the deep.”[14]
% mdparagraph:P118:end

\reserve{4\baselineskip}
% mdparagraph:P119:begin
\subsection*{Group I, Exercise 1}\BookHeading{nav-ch01-learner-0015}{subsubsection}{Group I, Exercise 1}{0}
% mdparagraph:P119:end

% item: EX-I-01-01
% mdparagraph:P120:begin
\itempara{1}{The evening and the morning.}
% mdparagraph:P120:end

% item: EX-I-01-02
% mdparagraph:P121:begin
\itempara{2}{The day and the night.}
% mdparagraph:P121:end

% item: EX-I-01-03
% mdparagraph:P122:begin
\itempara{3}{The earth and the heaven.}
% mdparagraph:P122:end

% item: EX-I-01-04
% mdparagraph:P123:begin
\itempara{4}{The fish and the bird.}
% mdparagraph:P123:end

% item: EX-I-01-05
% mdparagraph:P124:begin
\itempara{5}{The grass and the seed.}
% mdparagraph:P124:end

% item: EX-I-01-06
% mdparagraph:P125:begin
\itempara{6}{The tree and the herb.}
% mdparagraph:P125:end

% item: EX-I-01-07
% mdparagraph:P126:begin
\itempara{7}{The dry land and the water.}
% mdparagraph:P126:end

% item: EX-I-01-08
% mdparagraph:P127:begin
\itempara{8}{The image and the likeness.}
% mdparagraph:P127:end

% item: EX-I-01-09
% mdparagraph:P128:begin
\itempara{9}{The light and the darkness (plur.).}
% mdparagraph:P128:end

% item: EX-I-01-10
% mdparagraph:P129:begin
\itempara{10}{The sign and the season.}
% mdparagraph:P129:end

% item: EX-I-01-11
% mdparagraph:P130:begin
\itempara{11}{The spirit and the life.}
% mdparagraph:P130:end

% item: EX-I-01-12
% mdparagraph:P131:begin
\itempara{12}{The creature and the food.}
% mdparagraph:P131:end

% item: EX-I-01-13
% mdparagraph:P132:begin
\itempara{13}{The year and the day.}
% mdparagraph:P132:end

% item: EX-I-01-14
% mdparagraph:P133:begin
\itempara{14}{The man and the animal.}
% mdparagraph:P133:end

% item: EX-I-01-15
% mdparagraph:P134:begin
\itempara{15}{The tree and the fruit.}
% mdparagraph:P134:end

% item: EX-I-01-16
% mdparagraph:P135:begin
\itempara{16}{The firmament and the abyss.}
% mdparagraph:P135:end

\reserve{4\baselineskip}
% mdparagraph:P136:begin
\subsection*{Group I, Exercise 2}\BookHeading{nav-ch01-learner-0016}{subsubsection}{Group I, Exercise 2}{0}
% mdparagraph:P136:end

% item: EX-I-02-01
% mdparagraph:P137:begin
\itempara{1}{The days and the nights.}
% mdparagraph:P137:end

% item: EX-I-02-02
% mdparagraph:P138:begin
\itempara{2}{The fishes and the birds.}
% mdparagraph:P138:end

% item: EX-I-02-03
% mdparagraph:P139:begin
\itempara{3}{The trees and the herbs.}
% mdparagraph:P139:end

% item: EX-I-02-04
% mdparagraph:P140:begin
\itempara{4}{The images and the likenesses.}
% mdparagraph:P140:end

% item: EX-I-02-05
% mdparagraph:P141:begin
\itempara{5}{The signs and the seasons.}
% mdparagraph:P141:end

% item: EX-I-02-06
% mdparagraph:P142:begin
\itempara{6}{The years and the days.}
% mdparagraph:P142:end

% item: EX-I-02-07
% mdparagraph:P143:begin
\itempara{7}{The men and the animals.}
% mdparagraph:P143:end

% item: EX-I-02-08
% mdparagraph:P144:begin
\itempara{8}{The seas and the heavens.}
% mdparagraph:P144:end

% item: EX-I-02-09
% mdparagraph:P145:begin
\itempara{9}{The seeds and the fruits.}
% mdparagraph:P145:end

% item: EX-I-02-10
% mdparagraph:P146:begin
\itempara{10}{The evenings and the mornings.}
% mdparagraph:P146:end

% item: EX-I-02-11
% mdparagraph:P147:begin
\itempara{11}{The spirits and the men.}
% mdparagraph:P147:end

% item: EX-I-02-12
% mdparagraph:P148:begin
\itempara{12}{The whales and the cattle.}
% mdparagraph:P148:end

\reserve{4\baselineskip}
% mdparagraph:P149:begin
\subsection*{Group I, Exercise 3}\BookHeading{nav-ch01-learner-0017}{subsubsection}{Group I, Exercise 3}{0}
% mdparagraph:P149:end

% item: EX-I-03-01
% mdparagraph:P150:begin
\itempara{1}{The beginning of the day.}
% mdparagraph:P150:end

% item: EX-I-03-02
% mdparagraph:P151:begin
\itempara{2}{The spirit of God.}
% mdparagraph:P151:end

% item: EX-I-03-03
% mdparagraph:P152:begin
\itempara{3}{The light of the life.}
% mdparagraph:P152:end

% item: EX-I-03-04
% mdparagraph:P153:begin
\itempara{4}{The life of the bird.}
% mdparagraph:P153:end

% item: EX-I-03-05
% mdparagraph:P154:begin
\itempara{5}{The food of the fish.}
% mdparagraph:P154:end

% item: EX-I-03-06
% mdparagraph:P155:begin
\itempara{6}{The firmament of the heaven.}
% mdparagraph:P155:end

% item: EX-I-03-07
% mdparagraph:P156:begin
\itempara{7}{The beasts of the earth.}
% mdparagraph:P156:end

% item: EX-I-03-08
% mdparagraph:P157:begin
\itempara{8}{The face of the deep.}
% mdparagraph:P157:end

% item: EX-I-03-09
% mdparagraph:P158:begin
\itempara{9}{The congregations of the waters.}
% mdparagraph:P158:end

% item: EX-I-03-10
% mdparagraph:P159:begin
\itempara{10}{The waters of the sea.}
% mdparagraph:P159:end

% item: EX-I-03-11
% mdparagraph:P160:begin
\itempara{11}{The fishes of the sea.}
% mdparagraph:P160:end

% item: EX-I-03-12
% mdparagraph:P161:begin
\itempara{12}{The birds of the air.}
% mdparagraph:P161:end

% item: EX-I-03-13
% mdparagraph:P162:begin
\itempara{13}{The array of the heavens and of the earth.}
% mdparagraph:P162:end

% item: EX-I-03-14
% mdparagraph:P163:begin
\itempara{14}{The works of God.}
% mdparagraph:P163:end

% item: EX-I-03-15
% mdparagraph:P164:begin
\itempara{15}{The kinds of the animals.}
% mdparagraph:P164:end

% item: EX-I-03-16
% mdparagraph:P165:begin
\itempara{16}{The stars of the heaven.}
% mdparagraph:P165:end

% item: EX-I-03-17
% mdparagraph:P166:begin
\itempara{17}{The seeds of the grass.}
% mdparagraph:P166:end

% item: EX-I-03-18
% mdparagraph:P167:begin
\itempara{18}{The illumination of the nights.}
% mdparagraph:P167:end

\reserve{4\baselineskip}
% mdparagraph:P168:begin
\section*{14. Composition}\BookHeading{nav-ch01-learner-0018}{subsection}{14. Composition}{1}
% mdparagraph:P168:end

% mdparagraph:P169:begin
Use the selected forms and constructions without supplied verb frames.
% mdparagraph:P169:end

% item: CH01-COMP01
% mdparagraph:P170:begin
\itempara{1}{The bird is in the tree.}
% mdparagraph:P170:end

% item: CH01-COMP02
% mdparagraph:P171:begin
\itempara{2}{The birds were in the trees.}
% mdparagraph:P171:end

% item: CH01-COMP03
% mdparagraph:P172:begin
\itempara{3}{The seed is the bird's food.}
% mdparagraph:P172:end

% item: CH01-COMP04
% mdparagraph:P173:begin
\itempara{4}{The fish is an animal.}
% mdparagraph:P173:end

% item: CH01-COMP05
% mdparagraph:P174:begin
\itempara{5}{The signs were in the sky.}
% mdparagraph:P174:end

% item: CH01-COMP06
% mdparagraph:P175:begin
\itempara{6}{The herbs are in the man's dwelling.}
% mdparagraph:P175:end

% item: CH01-COMP07
% mdparagraph:P176:begin
\itempara{7}{The light of the night was in the sky.}
% mdparagraph:P176:end

% item: CH01-COMP08
% mdparagraph:P177:begin
\itempara{8}{The animals' food is the fruit of the trees.}
% mdparagraph:P177:end

\reserve{4\baselineskip}
% mdparagraph:P178:begin
\section*{15. Short readings}\BookHeading{nav-ch01-learner-0019}{subsection}{15. Short readings}{1}
% mdparagraph:P178:end

% mdparagraph:P179:begin
These two connected passages are new composition using the vocabulary and grammar of this unit. Translate from the Old English before answering the questions.
% mdparagraph:P179:end

% item: CH01-MINI01
% mdparagraph:P180:begin
\itempara{1}{Where was the man, and where is he now? Identify the reference of each pronoun. Explain the number of the subject and predicate noun in the fourth sentence, and identify whose dwelling is mentioned at the close.}
% mdparagraph:P180:end

% mdparagraph:P181:begin
Se mann wæs on ðǣm sǣ. Hē is on ðǣm hofe. Ðā fiscas sindon on ðǣm hofe. Hīe sindon ðæs mannes męte. Ðæt hof is on ðǣre eorðan. Hit is ðæs mannes hof.
% mdparagraph:P181:end

% item: CH01-MINI02
% mdparagraph:P182:begin
\itempara{2}{Translate the passage and rewrite it in the present. Explain why \textbf{hēo} does not establish the creature's sex. Then rewrite just the first two sentences about several creatures in the same dwelling, in the present; preserve “animals” as the predicate noun. Identify the reference of \textbf{hit} in the original passage and the case of \textbf{gesceafte}.}
% mdparagraph:P182:end

% mdparagraph:P183:begin
Sēo gesceaft wæs on ðǣm hofe. Hēo wæs dēor. Ðæt sǣd wæs on ðǣm hofe. Hit wæs ðǣre gesceafte męte.
% mdparagraph:P183:end

\reserve{4\baselineskip}
% mdparagraph:P184:begin
\section*{16. The first day of Creation}\BookHeading{nav-ch01-learner-0020}{subsection}{16. The first day of Creation}{1}
% mdparagraph:P184:end

% mdparagraph:P185:begin
Cook, \emph{A First Book in Old English}, Reader I, pp.123–124. The following is a continuous extract through the close of the first day. Only note cues and printed line/\allowbreak{}page breaks have been removed.[15]
% mdparagraph:P185:end

% source: FB-I-CREATION-FIRST-DAY-01
% mdparagraph:P186:begin
\readingtext{On anginne gescēop God heofonan and eorðan. Sēo eorðe sōðlice wæs īdel and ǣmtig; and ðīestru wǣron ofer ðǣre nēowolnesse brādnesse; and Godes gāst wæs gefęred ofer wæteru. God cwæð ðā, “Geweorðe lēoht”; and lēoht wearð geworht. God geseah ðā ðæt hit gōd wæs; and hē gedǣlde ðæt lēoht fram ðǣm ðīestrum. And hēt ðæt lēoht Dæg, and ðā ðīestru Niht. Ðā wæs geworden ǣfen and morgen ān dæg.}
% mdparagraph:P186:end

\reserve{4\baselineskip}
% mdparagraph:P187:begin
\subsection*{Supplied help}\BookHeading{nav-ch01-learner-0021}{subsubsection}{Supplied help}{0}
% mdparagraph:P187:end

% teaching: CH01-READ-SUPPORT CH01-GRAM-ORDER
% mdparagraph:P188:begin
Use this help for the questions and translation. It supplies the actual forms and constructions occurring here; their independent production is not required.
% mdparagraph:P188:end

% mdtable:T20:begin
% course-table-style: 2026-10
% course-panel:T20-P01:begin
\par\begingroup
\fontsize{10.5}{13.2}\selectfont
\setlength{\LTleft}{0pt}\setlength{\LTright}{\fill}
\begin{longtable}{@{}>{\raggedright\arraybackslash}p{32.28mm}>{\raggedright\arraybackslash}p{88.908mm}@{}}
\toprule
\textbf{Expression} & \textbf{Local meaning and analysis} \\
\midrule
\endfirsthead
\toprule
\textbf{Expression} & \textbf{Local meaning and analysis} \\
\midrule
\endhead
\midrule
\multicolumn{2}{r}{\footnotesize Continued on the following page.}\\
\endfoot
\bottomrule
\endlastfoot
\textbf{On anginne} & “In the beginning”: a supplied temporal use of \textbf{on}, with dative \textbf{anginne}. \\
\textbf{gescēop} & “Created,” third-person singular preterite. In the opening, the subject \textbf{God} follows the verb; word order alone does not identify the object. \\
\textbf{heofonan} & “Heaven”: accusative \textbf{singular} of weak feminine \textbf{heofone}, a variant of the masculine \textbf{heofon} learned above. It is not masculine \textbf{heofon}'s plural. \\
\textbf{sōðlice; īdel and ǣmtig} & Connective “now/\allowbreak{}for its part”; “empty and void.” The two adjectives are predicate feminine nominative singular with \textbf{eorðe}. \\
\textbf{ofer ðǣre nēowolnesse brādnesse} & \textbf{Ofer} means “over” and here requires a dative head. Apply the nested-genitive construction of §9 to identify the head and its dependent. \\
\textbf{wæs gefęred ofer wæteru} & “Was borne/\allowbreak{}carried over the waters,” contextually “moved over the waters.” \textbf{Wæs gefęred} is a passive phrase; \textbf{gefęred} is a past participle. Here \textbf{ofer} has accusative plural \textbf{wæteru}, unlike the preceding dative construction. \\
\textbf{cwæð ðā} & “Then said”: \textbf{cwæð} is third-person singular preterite; this \textbf{ðā} is the adverb “then,” not the article. \\
\textbf{Geweorðe lēoht} & “Let light come to be.” \textbf{Geweorðe} is a finite third-person singular present subjunctive (also called optative), expressing a wish or command. \\
\textbf{lēoht wearð geworht} & “Light was made/\allowbreak{}came into being.” \textbf{Wearð} is a third-person singular preterite used here with past participle \textbf{geworht} in a passive phrase. \\
\textbf{geseah ðā ðæt hit gōd wæs} & “Then saw that it was good.” \textbf{Geseah} is third-person singular preterite; \textbf{ðā} is “then”; \textbf{ðæt} is the conjunction “that.” \textbf{Gōd} is a neuter nominative singular predicate adjective. The copula stands at the end of the embedded clause. Use the context and §8 for the reference of \textbf{hit}. \\
\textbf{hē gedǣlde} & \textbf{Gedǣlde} is “separated,” third-person singular preterite. Identify the reference of the subject pronoun from the preceding clause. The object and \textbf{fram} phrase use §9. \\
\textbf{hēt ðæt lēoht Dæg, and ðā ðīestru Niht} & “[He] called the light Day, and the darkness Night.” \textbf{Hēt} is third-person singular preterite, with God/\allowbreak{}he understood as subject. \textbf{Ðæt lēoht} and \textbf{ðā ðīestru} are accusative objects; \textbf{Dæg} and \textbf{Niht} are accusative names applied to them. Here \textbf{ðæt} and \textbf{ðā} are articles. \\
\textbf{Ðā wæs geworden ǣfen and morgen ān dæg.} & “Then there came to be evening and morning: one day.” \textbf{Wæs geworden} expresses a resulting state or completed coming-to-be, “had come to be”; \textbf{geworden} is a past participle. \textbf{Ān dæg} is “one day,” not the ordinal “first day.” Treat this closing formula as a whole. \\
\end{longtable}
\par\endgroup
% course-panel:T20-P01:end
% mdtable:T20:end

\reserve{4\baselineskip}
% mdparagraph:P189:begin
\subsection*{Questions and translation}\BookHeading{nav-ch01-learner-0022}{subsubsection}{Questions and translation}{0}
% mdparagraph:P189:end

% mdparagraph:P190:begin
The unfamiliar verb constructions remain supplied. Questions 4 and 7 check that help; question 8 combines it with the noun, pronoun, and copula rules learned in this unit. They do not require independent recall of the subjunctive, passive, or naming construction.
% mdparagraph:P190:end

% item: CH01-READ01
% mdparagraph:P191:begin
\itempara{1}{Using the local support, identify the subject and two objects of the first sentence. Is heofonan plural? Translate the sentence.}
% mdparagraph:P191:end

% item: CH01-READ02
% mdparagraph:P192:begin
\itempara{2}{Describe the earth and the darkness. Explain wǣron, then identify the head and genitive dependent in the phrase after ofer.}
% mdparagraph:P192:end

% item: CH01-READ03
% mdparagraph:P193:begin
\itempara{3}{With the supplied passive phrase, state what God's spirit was doing and over what. Explain why wæteru cannot be the selected dative plural form.}
% mdparagraph:P193:end

% item: CH01-READ04
% mdparagraph:P194:begin
\itempara{4}{Does God's quoted speech report an existing fact or command/\allowbreak{}wish for something to come about? State what follows.}
% mdparagraph:P194:end

% item: CH01-READ05
% mdparagraph:P195:begin
\itempara{5}{In God geseah ðā ðæt hit gōd wæs, use the supplied functions of ðā and ðæt, and identify the reference of hit from the context. Translate the clause.}
% mdparagraph:P195:end

% item: CH01-READ06
% mdparagraph:P196:begin
\itempara{6}{What did God separate, and from what? Identify the direct-object and separation phrases and their cases.}
% mdparagraph:P196:end

% item: CH01-READ07
% mdparagraph:P197:begin
\itempara{7}{Using the naming and closing-formula support, state the names given, the understood subject of hēt, and the number of days in the close.}
% mdparagraph:P197:end

% item: CH01-READ08
% mdparagraph:P198:begin
\itempara{8}{Translate the entire first-day passage from the printed Old English, using the supplied help. Preserve the sequence, command, naming and one-day close.}
% mdparagraph:P198:end

\reserve{4\baselineskip}
% mdparagraph:P199:begin
\section*{17. Review}\BookHeading{nav-ch01-learner-0023}{subsection}{17. Review}{1}
% mdparagraph:P199:end

% item: CH01-REVIEW01
% mdparagraph:P200:begin
\itempara{1}{A learner writes *dægas and *hwælas for the plurals of dæg and hwæl, then changes wæstm to *wastmas. Correct the intended plurals and explain which vowel alternation applies.}
% mdparagraph:P200:end

% item: CH01-REVIEW02
% mdparagraph:P201:begin
\itempara{2}{In \textbf{on ðāra fugla trēowum}, a learner takes \textbf{ðāra} to agree with \textbf{trēowum}. Translate the phrase; identify the head and dependent, giving the case and number of each. Explain the article's actual agreement.}
% mdparagraph:P201:end

% item: CH01-REVIEW03
% mdparagraph:P202:begin
\itempara{3}{List the possible case/\allowbreak{}number values of \textbf{steorran} in the selected paradigm. Which applies in \textbf{ðæs steorran lēoht}, and which in \textbf{Ðā steorran wǣron on ðǣm heofone}? Translate both. Explain what the context adds that the ending alone cannot tell you.}
% mdparagraph:P202:end

% item: CH01-REVIEW04
% mdparagraph:P203:begin
\itempara{4}{Correct the agreement in \textbf{Ðā ðīestru is on ðǣm heofone. Ðæt lēoht wǣron on ðǣm hofe.} Keep the first sentence present and the second past; explain the corrections. Then complete \textbf{Se mann gedǣlde \_\_\_\_\_ fram \_\_\_\_\_.} to mean “The man separated the birds from the animal.” Use \textbf{dēor} for “animal.” The supplied \textbf{gedǣlde} means “separated,” third-person singular preterite. Identify the case and number of both completed phrases.}
% mdparagraph:P203:end

\reserve{4\baselineskip}
% mdparagraph:P204:begin
\section*{Source notes}\BookHeading{nav-ch01-learner-0024}{subsection}{Source notes}{1}
% mdparagraph:P204:end

% mdparagraph:P205:begin
\itempara{1}{Cook, \emph{A First Book in Old English}, second edition (1895), §§2–4, pp.5–6 /\allowbreak{} PDF21–22. The orthography follows the selected edition; this is not a historical account of all spellings or a complete pronunciation lesson.}
% mdparagraph:P205:end

% mdparagraph:P206:begin
\itempara{2}{Cook, §§38–42, pp.26–27 /\allowbreak{} PDF42–43. Tables composed from the checked paradigms are teaching models; they do not claim separate textual attestation of every form.}
% mdparagraph:P206:end

% mdparagraph:P207:begin
\itempara{3}{Cook, §84, p.50 /\allowbreak{} PDF66. The masculine article is \textbf{se}; the source distinguishes emphatic \textbf{sē}. Instrumental syntax is reserved for later treatment.}
% mdparagraph:P207:end

% mdparagraph:P208:begin
\itempara{4}{Cook, §§43–44, pp.28–30 /\allowbreak{} PDF44–46; especially §43.2 for \textbf{dæg/\allowbreak{}hwæl} and §43.4 for syncope and retained vowels.}
% mdparagraph:P208:end

% mdparagraph:P209:begin
\itempara{5}{Cook, §47, pp.32–33 /\allowbreak{} PDF48–49. The tables retain the selected \textbf{trēowu, tācnu, wæteru, nīetenu, ǣfenu} patterns. The neuter \textbf{ǣfen} plural is derived from its checked pattern, not quoted from a newly located passage.}
% mdparagraph:P209:end

% mdparagraph:P210:begin
\itempara{6}{Cook, §51, pp.34–36 /\allowbreak{} PDF50–52; §51.1 and the vocabulary's \textbf{wyrt} reference, p.324 /\allowbreak{} PDF340, support the endingless accusatives. The alternative N/\allowbreak{}A plural \textbf{-e} and singular \textbf{-a} of the \textbf{-ung} group are source variants.}
% mdparagraph:P210:end

% mdparagraph:P211:begin
\itempara{7}{Cook, §53, p.37 /\allowbreak{} PDF53.}
% mdparagraph:P211:end

% mdparagraph:P212:begin
\itempara{8}{Cook, §§46,52, pp.31–32,36 /\allowbreak{} PDF47–48,52; \textbf{mann/\allowbreak{}męnn}, §17, p.13 /\allowbreak{} PDF29; \textbf{niht}, Reader p.125 with note 28 /\allowbreak{} PDF141; \textbf{sǣ}, vocabulary p.324 /\allowbreak{} PDF340. Wright and Wright, \emph{Old English Grammar}, second edition (1914), §388, p.187 /\allowbreak{} PDF207, explicitly marks genitive plural \textbf{*sǣwa} as reconstructed. It is not tested here.}
% mdparagraph:P212:end

% mdparagraph:P213:begin
\itempara{9}{Cook, §138, pp.78–79 /\allowbreak{} PDF94–95. The bounded indicative series is not the complete \textbf{wesan/\allowbreak{}bēon} system.}
% mdparagraph:P213:end

% mdparagraph:P214:begin
\itempara{10}{Cook, §§149–150, p.88 /\allowbreak{} PDF104; §189, p.102 /\allowbreak{} PDF118.}
% mdparagraph:P214:end

% mdparagraph:P215:begin
\itempara{11}{Cook, \emph{Exercises in Old English} (1895), Exercise 3 introductory note, p.2 /\allowbreak{} PDF10; \emph{First Book}, §153, pp.89–90 /\allowbreak{} PDF105–106, especially §153.i.}
% mdparagraph:P215:end

% mdparagraph:P216:begin
\itempara{12}{Cook, \emph{First Book}, §§160,164(a), pp.93–94 /\allowbreak{} PDF109–110; §166, p.95 /\allowbreak{} PDF111; §§168,172, pp.96–97 /\allowbreak{} PDF112–113. \textbf{Sealde} is printed in §117, p.72 /\allowbreak{} PDF88. The other supplied frames occur in the assigned Creation extract. Examples in §9 are new composition. Supplied finite forms do not introduce their full conjugations.}
% mdparagraph:P216:end

% mdparagraph:P217:begin
\itempara{13}{The noun choices follow Cook's English–Old English vocabulary and the checked \textbf{First Book} paradigms. \textbf{Wæstm}, with short \textbf{æ}, corrects Cook's \textbf{wǣstm}; this correction applies also to its inflected forms. The source spelling is preserved in the editorial record. \textbf{Gærs}, \textbf{męte}, and the selected noun genders are retained. \textbf{Dēor} is not restricted to modern English “deer”; \textbf{nīeten} supplies “cattle” in the plural.}
% mdparagraph:P217:end

% mdparagraph:P218:begin
\itempara{14}{All 46 prompts are reproduced from Cook, \emph{Exercises}, Group I, Exercises 1–3, pp.1–2 /\allowbreak{} PDF9–10. Their English wording and original numbering are unchanged. Directions and answers are new teaching material.}
% mdparagraph:P218:end

% mdparagraph:P219:begin
\itempara{15}{Cook, \emph{First Book}, Reader I, pp.123–124 /\allowbreak{} PDF139–140: the five prose lines on p.123 are editorially counted; the first three on p.124 are counted from the printed marginal line number 5. The extract preserves \textbf{heofonan, gefęred, Geweorðe}, and unlengthened \textbf{geseah}. Its apparatus and original lineation remain separately transcribed. All PDF page numbers in these notes are one-based and refer to the selected scans.}
% mdparagraph:P219:end

% mdparagraph:P220:begin
\itempara{16}{Cook, §81, pp.48–49 /\allowbreak{} PDF64–65, gives the personal-pronoun paradigms. The subject inventory selects \textbf{hēo} and \textbf{hīe}, with \textbf{ðū} following this unit's spelling convention. Ringe and Taylor, \emph{The Development of Old English} (2014), §1.1.2, p.6 /\allowbreak{} preview PDF21, explicitly states that the dual pronouns take plural verbs. Fulk's Genesis 19 reading, p.38 /\allowbreak{} PDF50, uses ordinary second- and first-person plurals for the two angels, illustrating that plurals can include two. Only nominative subject use is introduced here.}
% mdparagraph:P220:end

% mdparagraph:P221:begin
\itempara{17}{Fulk, \emph{An Introductory Grammar of Old English} (2020), glossary p.265 /\allowbreak{} PDF277, prints unpointed \textbf{cyn(n)}; §18, p.12 /\allowbreak{} PDF24, gives its \textbf{k} value. The stopped initial of \textbf{gærs} is a source-based inference: Wright and Wright, \emph{Old English Grammar}, second edition (1914), §§314–315, p.156 /\allowbreak{} PDF176, classify \textbf{græs} with guttural \textbf{g} and describe its stopped initial outcome; §§66 note 4 and 345, pp.43,171 /\allowbreak{} PDF63,191, connect \textbf{gærs} with late r-metathesis. Compare Wright §85 note 1, p.47 /\allowbreak{} PDF67, for breaking before this metathesis, and Fulk §72, p.46 /\allowbreak{} PDF58, for palatalization before breaking. The learner needs only the supplied sound, not this historical derivation.}
% mdparagraph:P221:end


% book-source-body:end
\par\label{end-ch01-learner}
\clearpage\endgroup
